\section{Petrine Initiation in the “Old World” of Pure Intellect}

The debate with the old workman continues, and he assures Peter he is no demagogue:

\begin{quotex}
I am not such a fool in my old age that, understanding what is true, I should deny it for the favor of the rabble.
\end{quotex}


Niceta proposes that Peter be the master of the debate, which should be between himself and the old man.

\begin{quotex}
But after we have disputed in the Grecian manner, and we have come to that point where no issue appears, then he himself, as filled with the knowledge of Elohim, shall openly and clearly disclose to us the truth on all matters, so that not we only, but also all who are around us as hearers shall learn the way of truth. And therefore now let him sit as umpire; and when either of us shall yield, then let him, taking up the matter, give an unquestionable judgment.
\end{quotex}


Niceta and Aquila know the doctrines of Epicurus and Pyrhonnists, but the old man says his doctrine is different: he follows ``Genesis''.

\begin{quotex}
I Clement said to him:Hear, my father: if my brother Niceta bring you to acknowledgement that the world is not governed without the providence of Elohim, I shall be able to answer you in that part which remains concerning \emph{the Genesis}; for I am well acquainted with this doctrine.
\end{quotex}


A discussion of simplicity or ``substance'' follows, which ends with a highlighting of mathematics:

Fora way is afforded us to intellectual and invisible things from those which we see and handle; as is contained in arithmetical instructions, where, when inquiry is made concerning divine things,we rise from the lower to the higher numbers; but when the method respecting present and visible things is expounded, the order is directed from the higher to the lower numbers. Is it not so?

``Genesis'' turns out to be roughly Deism, of some sorts; Niceta objects that prayer, then, is futile, and that on the contrary, Elohim will destroy the present ``clockwork'' world and reward those with ageless life who are worthy of inheriting the invisible world to come, which contains the visible world. We are on very Platonic ground -- the spirit is even said to \emph{escape} from the body.Matter is made of the four elements, tempered by a certain something of Elohim's action.

For I think that what some have said is vain: that the body of the world is simple, that is, without any conjunction; since it is evident that what is simple can neither be a body, nor can be mixed,or propagated, or dissolved; all of which, we see, the bodies of the world do.

A learned discussion follows, summarizing Greek philosophy:

\begin{quotex}
For the Greek philosophers, inquiring into the beginnings of the world, have gone some in one way and some in another. In short,Pythagoras says that numbers are the elements of its beginnings;Callistratus says qualities; Alcmaeon, contrarieties; Anaximander,immensity; Anaxagoras, equalities of parts; Epicurus, atoms;Diodorus, akatonomaston (that is, things in which there are no parts; Asclepius, Ogkoi,which we may call tumors or swellings; the geometricians, ends; Democritus, ideas; Thales, water;Heraclitus, fire; Diogenes, air; Parmenides, earth; Zeno, Empedocles,and Plato, fire, water, air, and earth. Aristotle also introduces a fifth element, which he called akatonomaston; that is, that which cannot be named; without doubt indicating Elohim who made the world, by joining the four elements into one.
\end{quotex}


Extremely good arguments are made against Epicurus, who is taken as an exemplar of the materialist schools -- they seem to make a teleological argument for God.

\begin{quotex}
and that insensible bodies,even if they could by any means concur and be united, could not give forms and measures to bodies, form limbs, or effect qualities, or express quantities; all which, therefore, by their exactness, attest the hand of a Maker, and show the operation of reason, which reason I call the Word, and YHWH
\end{quotex}


Plato, in the \emph{Timaeus}, is appealed to for belief in a Maker.

Apparently, the speakers shared the medieval and ancient beliefs in the changing of sex (eg., in hyenas) \& in spontaneous generation from dung. Again, water has a special place (which reminds me of Lao Tzu\footnote{\url{https://en.wikipedia.org/wiki/Tao\_Te\_Ching}}):

But to prove by facts and examples that nothing is imparted to seeds of the substance of the earth, but that all depends upon the element of water, and the power of the Ruach which is in it... The doctrine of the Green Tablet:

\begin{quotex}
But let us come now, if you please, to our own substance, that is, the substance of man, who is a small world, a microcosm,in the great world; and let us consider with what reason it is compounded: andfrom this especially you will understand the Hokmah of the Creator... But if it is Reason that is, Logos by which it appears that all things were made, they change the name without purpose, when they make statements concerning the reason of the Creator. If you have anything to say to these things, my father, say on.
\end{quotex}


In an interlude, Peter emphasizes the need for doctrinal faithfulness:

\begin{quotex}
yet so that you only join the eloquence of your discourse with those things which you have heard from me, and which have been committed to you. But do not speak anything which is your own, and which has not been committed to you, though it may seem to yourselves to be true; but hold forth those things, as I have said, which I myself have received from the Navi Emet, Yshua,and have delivered to you, although they may seem to be less full of authority. For thus men often do who turn away from the truth,while they believe that they have found out, by their own thoughts, a form of truth more true and powerful.
\end{quotex}


The old man basically is a modern sceptic who objects to the ``Watchmaker Fallacy''\footnote{\url{https://en.wikipedia.org/wiki/Watchmaker\_analogy}}:

Then said the old man:

\begin{quotex}
Behold, we see the bow in the heaven assume a circular shape,completed in all proportion, and have an appearance of reality,which in all probability neither mind could have constructed nor reason described; and yet it is not made by any human mind. Behold,I have set forth the whole in a word: now answer me.
\end{quotex}


Aquila counters with Augustine's doctrine of seminal Reason\footnote{\url{https://en.wikipedia.org/wiki/Rationes\_seminales}}:

\begin{quotex}
But this is similar to that, that in the beginning Elohim created insensible elements that He might use for forming and developing all other things. But even those who form statues first make a mold of clay or wax, and from it the figure of the statue is produced.
\end{quotex}


It quickly reduces to the Problem of Evil \& the Problem of Good:

\begin{quotex}
To this the old man added:I am not able, my son, to say that those things which seem formed by art are made by mind, by reason of other things which we see to be done unjustly and disorderly in the world. If, says Aquila, those things which are done disorderly do not allow you say that they are done by YHWHs providence, why dont those things which are done orderly compel you to say that they are done by YHWH, and that irrational nature cannot produce a rational work? For it is certain, nor do we at all deny, that in this world somethings are done orderly, and some disorderly.
\end{quotex}


Evil falls on all, but redounds to the advantage of the good:

\begin{quotex}
But for this reason YHWH has decreed a judgment with respect to all men, \emph{because the present life was not such that everyone could be dealt with according to his deserving.} Those things,therefore, which were well and orderly appointed from the beginning, when no causes of evil existed, are not to be judged of from the evils which have befallen the world by reason of the sins of men.
\end{quotex}


The Seres tribe\footnote{\url{https://en.wikipedia.org/wiki/Seres}} is mentioned as living ``without evil''.

Some form of natural religion is endorsed in the sense that Good \& Evil are always intertwined, \& man's business is to sort them:

\begin{quotex}
But do you, my father, as a wise man, choose from that division the part which preserves order and makes for the belief of providence, and do not only follow that part which runs against order and neutralizes the belief of Providence... .To this the old man answered: Show me a way, my son, by which I may establish in my mind one or the other of these two orders, the one of which asserts, and the other denies, providence.Aquila answered,To one having a right judgment the decision is easy. For this very thing that you say, order and disorder, may be produced by a contriver, but not by insensible nature. For let us suppose, by way of illustration, that a great mass were torn from a high rock, and cast down headlong, and when clashed upon the ground were broken into many pieces, could it in any way be that, amongst that multitude of fragments, there should be found even one which should have any perfect figure and shape?The old man answered:It is impossible.But,said Aquila, if there be present a statuary, he can by his skilful hand and reasonable mind form the stone cut from the mountain into whatever figure he pleases. The old man said: That is true. Therefore,says Aquila, when there is not a rational mind, no figure can be formed out of the mass; but when there is a designing mind, there may be both form and deformity: for example, if a workman cuts from the mountain a block to which he wishes to give a form, he must first cut it out unformed and rough; then, by degrees hammering and hewing it by the rule of his art, he expresses the form which he has conceived in his mind. Thus, therefore, from infirmity or deformity, by the hand of the workman form is attained,and both proceed from the workman. In like manner, therefore, the things which are done in the world are accomplished by the providence of a contriver, although they may seem not quite orderly. And therefore, because these two ways have been made known to you, and you have heard the divisions of them, flee from the way of unbelief, lest it lead you to that prince who delights in evils; but follow the way of faith, that you may come to that King who delights in good men.
\end{quotex}


This is truly Augustine's De Vera Religione\footnote{\url{http://www.scribd.com/doc/86329463/Augustine-on-True-Religion-De-Vera-Religione}}. Apparently, some form of predestination also explains the origin of evil -- men given to evil and evil rulers or demons were foreseen by God to have chosen them already. That is, the eternal state of individual being explains the transitory manifestation -- thus the problem of evil is set free from the normal Western modes of wrestling with it.

\begin{quotex}
And this is the bound assigned, that unless one first do the will of the demons, the demons have no power over him.
\end{quotex}


Thus, states of Being are self-chosen.

\emph{Deucalion\footnote{\url{https://en.wikipedia.org/wiki/Deucalion}} is mentioned in regards to the Flood (this is an old ``myth'').}

The old persists, pushing the argument back to origins:

\begin{quotex}
You have stated it excellently, my son. It now remains only that you tell me whence is the substance of evil.For if it was made by Elohim, the evil fruit shows that the root is at fault; for it appears that it also is of an evil nature. But if this substance was co-ageless with Elohim, how can that which was equally unproduced and co-ageless be subject to the other... 
\end{quotex}


As Boethius would argue a few centuries later, the weight of evidence falls on the power of the rational mind dominated by Logos, to turn the scales and decide this question (along with all others). Please note this argument holds no weight for the irrational and those dominated by passions:

\begin{quotex}
but if a reasonable mind, which has been made by Elohim, do not acquiesce in the Torah of its Creator, and go beyond the bounds of the temperance prescribed to it, how does this reflect on the Creator? Or if there is any reason higher than this, we do not know it; for we cannot know anything perfectly, and especially concerning those things for our ignorance of which we are not to be judged. But those things for which we are to be judged are most easy to be understood, and are dispatched almost in a word.For almost the whole rule of our actions is summed up in this, that what we are unwilling to suffer we should not do to others. For as you would not be killed, you must beware of killing another; and as you would not have your own marriage violated, you must not defile another's bed; you would not be stolen from, neither must you steal;and every matter of men's actions is comprehended within this rule.
\end{quotex}


The Golden Rule (or the First and Second Commandments) have Ontology and Epistemology behind them -- they are not sentimental -- one can think of Carlyle's Eternal No\footnote{\url{https://en.wikipedia.org/wiki/Thomas\_Carlyle}} -- man knows (at least) how to begin the journey and \emph{abstain from evil previously chosen}.

The old man is unconvinced, because he does not think it within man's power to ultimately choose good over evil as if free from Fate, \& Aquila offers to let Clement argue mathematically, since he is better educated in the wisdom of numbers. Simon Peter reminds all of them that they can know nothing of God unless it is revealed, for even with men, we do not know a man's mind until he speaks.

\begin{quotex}
How much more must it be that no one can know the mind or the work of the invisible and incomprehensible Elohim, unless He Himself sends a \emph{navi} to declare His purpose and expound the way of His creation, so far as it is lawful for men to learn it! Hence I think it ridiculous when men judge of the power of Elohim in natural ways, and think that this is possible and that impossible to Him, or this greater and that less,while they are ignorant of everything; who, being unrighteous, judge the righteous YHWH; unskilled, judge the contriver; corrupt, judge the incorruptible; created, judge the Creator... 
\end{quotex}

A Jewish argument, certainly, but one compatible with the tradition of the Emerald Tablet -- Man (the microcosm) needs a microcosm of the macrocosm to speak to him.Making it even more explicit, one wonders how ``Semiticism'' stayed Semitic -- this sounds positively Greek -- the eternal Logos:

And therefore I advise not only wise men, but indeed all men who have a desire of knowing what is advantageous to them, that they seek after the Navi Emet Yshua; for it is He alone who knoweth all things, and who knoweth what and how every man is seeking. \emph{For He is within the mind of every one of us, but in those who have no desire of the knowledge of Elohim and His zedekah, \textbf{He is inoperative;} but He works in those who seek after that which is profitable to their spirits, and kindles in them the light of knowledge}.Wherefore seek Him first of all; and if you do not find Him, do not expect that you shall learn anything from any other. But He is soon found by those who diligently seek Him through love of the truth,and whose spirits are not taken possession of by wickedness... \emph{For He is present with those who desire Him in the innocence of their spirits, who bear patiently, \textbf{and draw sighs from the bottom of their hearts} through love of the truth}.

What else is this but proof that Christianity originated as an esoteric religion (first and foremost) which was then made exoteric, and attempted to draw existing exoteric religions/practices/rites back into the bosom of its own interiority?

\begin{quotex}
Therefore, if any one wishes to learn all things, he must do so little by little, for, being mortal, he shall not be able to comprehend all at once the counsel of Elohim and to scan immensity itself. But if, as we have said, he desires to learn all things, let him seek after the Navi Emet; and when he has found Him, let him not treat with Him by questions and disputations and arguments; but if He has given any response, or pronounced any judgment, it cannot be doubted that this is certain. And therefore, before all things, let the Navi Emet be sought, and His words be laid hold of... \end{quotex}


The stony deserts of the wordy are warned against, explicitly, and an argument offered for remembering both Sweetness and Light\footnote{\url{https://en.wikipedia.org/wiki/Sweetness\_and\_light}}:

\begin{quotex}
And for the rest may they be at \emph{shalom}, having received evident knowledge of the truth. For all other things are treated by opinion, in which there can be nothing firm. For what speech is there which may not be contradicted? And what argument is there that may not be overthrown by another argument? And hence it is, that by disputation of this sort men can never come to any end of knowledge and learning, but find the end of their life sooner than the end of their questions... 
\end{quotex}


Hopefully, this gives a better insight into the origins of Christianity... the journal is only two-thirds over... .


\flrightit{Posted on 2012-08-24 by Logres}

